\FigureLayoutDeclare{img/2006/jiangsu/q18_fig1.png}{4.0cm}{2bp 2bp 2bp 2bp}
\FigureLayoutDeclare{img/2006/jiangsu/q9_fig1.png}{12.0cm}{22.91bp 0.655bp 23.237bp 4.909bp}

\chapter{2006年江苏卷}
\section{单选题}
\begin{problem}
已知 \(a \in \R\)，函数 \(f(x) = \sin x - |a|\)，\(x \in \R\) 为奇函数，则 \(a =\)（\quad）

\choices
    {\(0\)}
    {\(1\)}
    {\(-1\)}
    {\(\pm1\)}
\end{problem}

\begin{answer}
A
\end{answer}

\begin{problem}
圆 \((x - 1)^2 + (y + \sqrt{3})^2 = 1\) 的切线方程中有一个是（\quad）

\choices
    {\(x - y = 0\)}
    {\(x + y = 0\)}
    {\(x = 0\)}
    {\(y = 0\)}
\end{problem}

\begin{answer}
C
\end{answer}

\begin{problem}
某人 \(5\) 次上班途中所花的时间（单位：\text{分钟}）分别为 \(x\)，\(y\)，\(10\)，\(11\)，\(9\). 已知这组数据的平均数为 \(10\)，方差为 \(2\)，则 \(|x - y|\) 的值为（\quad）

\choices
    {\(1\)}
    {\(2\)}
    {\(3\)}
    {\(4\)}
\end{problem}

\begin{answer}
D
\end{answer}

\begin{problem}
为了得到函数 \(y = 2\sin \left(\frac{x}{3} + \frac{\pi}{6}\right)\)，\(x \in \R\) 的图象，只需把函数 \(y = 2\sin x\)，\(x \in \R\) 的图象上所有的点（\quad）

\choices
    {向左平移 \(\frac{\pi}{6}\) 个单位长度，再把所得各点的横坐标缩短到原来的 \(\frac{1}{3}\) 倍（纵坐标不变）}
    {向右平移 \(\frac{\pi}{6}\) 个单位长度，再把所得各点的横坐标缩短到原来的 \(\frac{1}{3}\) 倍（纵坐标不变）}
    {向左平移 \(\frac{\pi}{6}\) 个单位长度，再把所得各点的横坐标伸长到原来的 \(3\) 倍（纵坐标不变）}
    {向右平移 \(\frac{\pi}{6}\) 个单位长度，再把所得各点的横坐标伸长到原来的 \(3\) 倍（纵坐标不变）}
\end{problem}

\begin{answer}
C
\end{answer}

\begin{problem}
\(\left(\sqrt{x} -\frac{1}{3x}\right)^{10}\) 的展开式中含 \(x\) 的正整数指数幂的项数是（\quad）

\choices
    {\(0\)}
    {\(2\)}
    {\(4\)}
    {\(6\)}
\end{problem}

\begin{answer}
B
\end{answer}

\begin{problem}
已知两点 \(M(-2,0)\)，\(N(2,0)\)，点 \(P\) 为坐标平面内的动点，满足 \(\left|\overrightarrow{MN}\right| \cdot \left|\overrightarrow{MP}\right| + \overrightarrow{MN} \cdot \overrightarrow{NP} = 0\)，则动点 \(P(x,y)\) 的轨迹方程为（\quad）

\choices
    {\( y^{2} = 8x \)}
    {\( y^{2} = -8x \)}
    {\( y^{2} = 4x \)}
    {\( y^{2} = -4x \)}
\end{problem}

\begin{answer}
B
\end{answer}

\begin{problem}
若 \(A\)，\(B\)，\(C\) 为三个集合，\(A \cup B = B \cap C\)，则一定有（\quad）

\choices
    {\(A \subseteq C\)}
    {\(C \subseteq A\)}
    {\(A \neq C\)}
    {\(A \neq \emptyset\)}
\end{problem}

\begin{answer}
A
\end{answer}

\begin{problem}
设 \(a\)，\(b\)，\(c\) 是互不相等的正数，则下列不等式中不恒成立的是（\quad）

\choices
    {\( |a - b| \leqslant |a - c| + |b - c| \)}
    {\(a^2 +\frac{1}{a^2}\geqslant a + \frac{1}{a}\)}
    {\( |a - b| + \frac{1}{a - b} \geqslant 2 \)}
    {\(\sqrt{a + 3} -\sqrt{a + 1}\leqslant \sqrt{a + 2} -\sqrt{a}\)}
\end{problem}

\begin{answer}
C
\end{answer}

\begin{problem}
两个相同的正四棱锥组成如图所示的几何体，可放入棱长为 \(1\) 的正方体内，使正四棱锥的底面 \(ABCD\) 与正方体的某一个平面平行，且各顶点均在正方体的面上，则这样的几何体体积的可能值有（\quad）

\bitmapfigure[max width=0.78\linewidth]{img/2006/jiangsu/q9_fig1.png}

\choices
    {\(1\) 个}
    {\(2\) 个}
    {\(3\) 个}
    {无穷多个}
\end{problem}

\begin{answer}
D
\end{answer}

\begin{problem}
图中有一个信号源和五个接收器. 接收器与信号源在同一个串联线路中时，就能接收到信号，否则就不能接收到信号. 若将图中左端的六个接线点随机地平均分成三组. 将右端的六个接线点也随机地平均分成三组，再把所有六组中每组的两个接线点用导线连接，则这五个接收器能同时接收到信号的概率是（\quad）

\bitmapfigure[max width=0.42\linewidth]{img/2006/jiangsu/q10_fig1.png}

\choices
    {\(\frac{4}{45}\)}
    {\(\frac{1}{36}\)}
    {\(\frac{4}{15}\)}
    {\(\frac{8}{15}\)}
\end{problem}

\begin{answer}
D
\end{answer}

\section{填空题}
\begin{problem}
在 \(\triangle ABC\) 中，已知 \(BC = 12\)，\(A = 60^{\circ}\)，\(B = 45^{\circ}\)，则 \(AC = \)\fillinblank{}.
\end{problem}

\begin{answer}
\(4\sqrt{6}\)
\end{answer}

\begin{problem}
设变量 \(x\)，\(y\) 满足约束条件 \(\begin{cases}
2x - y \leqslant 2,\\
x - y \geqslant -1,\\
x + y \geqslant 1.
\end{cases}\)，则 \(z = 2x + 3y\) 的最大值为\fillinblank{}.
\end{problem}

\begin{answer}
\(18\)
\end{answer}

\begin{problem}
今有 \(2\) 个红球，\(3\) 个黄球，\(4\) 个白球，同色球不加以区分，将这 \(9\) 个球排成一列有\fillinblank{} 种不同的方法.（用数字作答）
\end{problem}

\begin{answer}
\(1260\)
\end{answer}

\begin{problem}
\(\cot 20^{\circ}\cos 10^{\circ} + \sqrt{3}\sin 10^{\circ}\tan 70^{\circ} - 2\cos 40^{\circ} =\)\fillinblank{}.
\end{problem}

\begin{answer}
\(2\)
\end{answer}

\begin{problem}
对正整数 \(n\)，设曲线 \(y = x^n (1 - x)\) 在 \(x = 2\) 处的切线与 \(y\) 轴交点的纵坐标为 \(a_n\)，则数列 \(\left\{\frac{a_n}{n + 1}\right\}\) 的前 \(n\) 项和的公式是\fillinblank{}.
\end{problem}

\begin{answer}
\(2^{n+1}-2\)
\end{answer}

\begin{problem}
不等式 \(\log_2\left(x + \frac{1}{x} + 6\right)\leqslant 3\) 的解集为\fillinblank{}.
\end{problem}

\begin{answer}
\(\left(-3-2\sqrt{2},-3+2\sqrt{2}\right)\cup\{1\}\)
\end{answer}

\section{解答题}
\begin{problem}
已知三点 \(P(5,2)\)，\(F_{1}(-6,0)\)，\(F_{2}(6,0)\).

\begin{enumerate}
\item 求以 \(F_{1}\)，\(F_{2}\) 为焦点且过点 \(P\) 的椭圆的标准方程；
\item 设点 \(P\)，\(F_{1}\)，\(F_{2}\) 关于直线 \(y = x\) 的对称点分别为 \(P'\)，\(F_{1}'\)，\(F_{2}'\)，求以 \(F_{1}'\)，\(F_{2}'\) 为焦点且过点 \(P'\) 的双曲线的标准方程.
\end{enumerate}
\end{problem}

\begin{solution}
\begin{enumerate}
\item 设所求椭圆的标准方程为 \(\dfrac{x^2}{a^2}+\dfrac{y^2}{b^2}=1\)，其中 \(a>b>0\)，焦距参数为 \(c\). 由焦点坐标知 \(c=6\)，且
\(PF_1=\sqrt{(5+6)^2+2^2}=5\sqrt{5}\)，\(PF_2=\sqrt{(5-6)^2+2^2}=\sqrt{5}\). 因为椭圆上任一点到两焦点的距离之和等于 \(2a\)，所以 \(2a=6\sqrt{5}\)，即 \(a^2=45\). 又 \(b^2=a^2-c^2=45-36=9\)，故椭圆的标准方程为 \(\dfrac{x^2}{45}+\dfrac{y^2}{9}=1\).
\item 关于直线 \(y=x\) 对称后，\(P'=(2,5)\)，\(F_1'=(0,-6)\)，\(F_2'=(0,6)\). 设所求双曲线的标准方程为 \(\dfrac{y^2}{a_1^2}-\dfrac{x^2}{b_1^2}=1\)，则焦距参数仍为 \(c_1=6\). 由双曲线定义，\(2a_1=|P'F_1'-P'F_2'|=\sqrt{2^2+11^2}-\sqrt{2^2+(-1)^2}=4\sqrt{5}\)，所以 \(a_1^2=20\). 因而 \(b_1^2=c_1^2-a_1^2=36-20=16\)，所求双曲线的标准方程为 \(\dfrac{y^2}{20}-\dfrac{x^2}{16}=1\).
\end{enumerate}
\end{solution}

\begin{problem}
请您设计一个帐篷. 它下部的形状是高为 \(1\mathrm{~m}\) 的正六棱柱，上部的形状是侧棱长为 \(3\mathrm{~m}\) 的正六棱锥（如图所示）. 试问当帐篷的顶点 \(O\) 到底面中心 \(O_{1}\) 的距离为多少时，帐篷的体积最大？

\bitmapfigure[max width=0.42\linewidth]{img/2006/jiangsu/q18_fig1.png}
\end{problem}

\begin{solution}
设 \(OO_{1}=x\mathrm{~m}\). 由于正六棱柱的高为 \(1\mathrm{~m}\)，所以正六棱锥的高为 \((x-1)\mathrm{~m}\). 设正六棱锥底面正六边形的边长为 \(s\)，正六边形的中心到顶点的距离等于边长，故由侧棱长为 \(3\mathrm{~m}\) 得
\(s^2+(x-1)^2=3^2\)，即 \(s^2=8+2x-x^2\). 正六棱锥非退化时 \(1<x<4\).

底面正六边形的面积为 \(S=\dfrac{3\sqrt{3}}{2}s^2=\dfrac{3\sqrt{3}}{2}(8+2x-x^2)\)，于是帐篷的体积为
\(V(x)=S\cdot1+\dfrac13S(x-1)=\dfrac{\sqrt{3}}{2}(8+2x-x^2)(x+2)\).
因此 \(V'(x)=\dfrac{\sqrt{3}}{2}(12-3x^2)\). 当 \(1<x<2\) 时，\(V'(x)>0\)；当 \(2<x<4\) 时，\(V'(x)<0\). 所以当 \(x=2\) 时体积最大，即 \(OO_{1}=2\mathrm{~m}\).
\end{solution}

\begin{problem}
在正三角形 \(ABC\) 中，\(E\)，\(F\)，\(P\) 分别是 \(AB\)，\(AC\)，\(BC\) 边上的点. 满足 \(AE:EB=CF:FA=CP:PB=1:2\)（如图1）. 将 \(\triangle AEF\) 沿 \(EF\) 折起到 \(\triangle A_{1}EF\) 的位置. 使二面角 \(A_{1}-EF-B\) 成直二面角，连接 \(A_{1}B\)，\(A_{1}P\)（如图2）.

\begin{enumerate}
\item 求证：\(A_{1}E \perp\) 平面 \(BEP\)；
\item 求直线 \(A_{1}E\) 与平面 \(A_{1}BP\) 所成角的大小；
    \item 求二面角 \(B - A_{1}P - F\) 的大小.（用反三角函数表示）
\end{enumerate}

\bitmapfigure[max width=0.54\linewidth]{img/2006/jiangsu/q19_fig1.png}
\end{problem}

\begin{solution}
\begin{enumerate}
\item 取正三角形 \(ABC\) 的边长为 \(3\)，建立空间直角坐标系，使折叠前的平面 \(ABC\) 为 \(z=0\)，并取
\(B=(0,0,0)\)，\(C=(3,0,0)\)，\(A=\left(\dfrac32,\dfrac{3\sqrt3}{2},0\right)\). 由题设的分点关系可得 \(E=(1,\sqrt3,0)\)，\(F=\left(\dfrac52,\dfrac{\sqrt3}{2},0\right)\)，\(P=(2,0,0)\).

由 \(AE=1\)，\(AF=2\)，\(EF=\sqrt3\)，有 \(AE^2+EF^2=AF^2\)，所以 \(AE\perp EF\). 折叠后 \(A_1E\perp EF\) 仍成立；又因为二面角 \(A_1-EF-B\) 为直二面角，平面 \(A_1EF\perp\) 平面 \(BEF\)，所以平面 \(A_1EF\) 内垂直于 \(EF\) 的直线 \(A_1E\) 垂直于平面 \(BEF\). 由于 \(B,E,P\) 均在原平面 \(ABC\) 内，平面 \(BEP=BEF\)，故 \(A_1E\perp\) 平面 \(BEP\).
\item 由 \(A_1E\perp\) 平面 \(ABC\) 且 \(A_1E=AE=1\)，取图示折叠方向可得 \(A_1=(1,\sqrt3,1)\). 在平面 \(A_1BP\) 内取 \(\overrightarrow{BP}=(2,0,0)\)，\(\overrightarrow{BA_1}=(1,\sqrt3,1)\)，则该平面的一个法向量为
\(\bs{n}=\overrightarrow{BP}\times\overrightarrow{BA_1}=(0,-2,2\sqrt3)\). 取直线 \(A_1E\) 的方向向量 \(\bs{u}=(0,0,1)\)，设直线与平面所成的角为 \(\theta\)，则
\(\sin\theta=\dfrac{|\bs{u}\cdot\bs{n}|}{|\bs{u}|\,|\bs{n}|}=\dfrac{2\sqrt3}{4}=\dfrac{\sqrt3}{2}\). 因为 \(0<\theta<\dfrac\pi2\)，所以 \(\theta=60^{\circ}\).
\item 取过 \(P\) 且垂直于棱 \(A_1P\) 的平面作截面. 令 \(\bs{u}=\overrightarrow{PA_1}=(-1,\sqrt3,1)\)，则 \(|\bs{u}|^2=5\). 记 \(\bs{b}=\overrightarrow{PB}=(-2,0,0)\)，\(\bs{f}=\overrightarrow{PF}=\left(\dfrac12,\dfrac{\sqrt3}{2},0\right)\)，它们在该截面上的投影分别为
\(\bs{b}_{\perp}=\bs{b}-\dfrac{\bs{b}\cdot\bs{u}}{|\bs{u}|^2}\bs{u}=\left(-\dfrac85,-\dfrac{2\sqrt3}{5},-\dfrac25\right)\)，
\(\bs{f}_{\perp}=\bs{f}-\dfrac{\bs{f}\cdot\bs{u}}{|\bs{u}|^2}\bs{u}=\left(\dfrac7{10},\dfrac{3\sqrt3}{10},-\dfrac15\right)\).

于是 \(\bs{b}_{\perp}\cdot\bs{f}_{\perp}=-\dfrac75\)，\(|\bs{b}_{\perp}|=\dfrac4{\sqrt5}\)，\(|\bs{f}_{\perp}|=\dfrac2{\sqrt5}\). 设二面角 \(B-A_1P-F\) 为 \(\varphi\)，则截面内两射线的夹角满足
\(\cos\varphi=\dfrac{\bs{b}_{\perp}\cdot\bs{f}_{\perp}}{|\bs{b}_{\perp}|\,|\bs{f}_{\perp}|}=-\dfrac78\). 因此 \(\varphi=\arccos\left(-\dfrac78\right)=\pi-\arccos\dfrac78\).
\end{enumerate}
\end{solution}

\begin{problem}
设 \(a\) 为实数，设函数 \(f(x) = a\sqrt{1 - x^2} + \sqrt{1 + x} + \sqrt{1 - x}\) 的最大值为 \(g(a)\).

\begin{enumerate}
\item 设 \(t = \sqrt{1 + x} + \sqrt{1 - x}\)，求 \(t\) 的取值范围，并把 \(f(x)\) 表示为 \(t\) 的函数 \(m(t)\)；
\item 求 \(g(a)\)；
\item 试求满足 \(g(a) = g\left(\frac{1}{a}\right)\) 的所有实数 \(a\).
\end{enumerate}
\end{problem}

\begin{solution}
\begin{enumerate}
\item 由根式有意义知 \(-1\leqslant x\leqslant1\). 令 \(t=\sqrt{1+x}+\sqrt{1-x}\)，则
\(t^2=2+2\sqrt{1-x^2}\)，所以 \(\sqrt2\leqslant t\leqslant2\). 由 \(\sqrt{1-x^2}=\dfrac{t^2-2}{2}\)，得
\(f(x)=a\sqrt{1-x^2}+t=\dfrac{a}{2}t^2+t-a=:m(t)\)，其中 \(t\in[\sqrt2,2]\).
\item
\examdisplaycases{g(a) =}{
a+2, & a>-\dfrac{1}{2},\\
-a-\dfrac{1}{2a}, & -\dfrac{\sqrt{2}}{2}<a\leqslant-\dfrac{1}{2},\\
\sqrt{2}, & a\leqslant-\dfrac{\sqrt{2}}{2}.
}
由 \(m'(t)=at+1\) 可分类讨论. 当 \(a>-\dfrac12\) 时，函数在区间 \([\sqrt2,2]\) 上的最大值在 \(t=2\) 处取得，故 \(g(a)=m(2)=a+2\). 当 \(-\dfrac{\sqrt2}{2}<a\leqslant-\dfrac12\) 时，顶点 \(t_0=-\dfrac1a\in[\sqrt2,2]\)，故 \(g(a)=m(t_0)=-a-\dfrac1{2a}\). 当 \(a\leqslant-\dfrac{\sqrt2}{2}\) 时，\(m(t)\) 在该区间上递减，故 \(g(a)=m(\sqrt2)=\sqrt2\).
\item 因为 \(g\left(\dfrac1a\right)\) 有意义，所以先要求 \(a\neq0\). 当 \(a>0\) 时，\(g(a)=a+2\)，\(g\left(\dfrac1a\right)=\dfrac1a+2\)，于是 \(a=1\). 当 \(a<0\) 且 \(-\sqrt2\leqslant a\leqslant-\dfrac{\sqrt2}{2}\) 时，\(\dfrac1a\) 仍在区间 \(\left[-\sqrt2,-\dfrac{\sqrt2}{2}\right]\) 内，所以两边都等于 \(\sqrt2\). 若 \(a<-\sqrt2\)，则 \(g(a)=\sqrt2\)，而 \(\dfrac1a>-\dfrac{\sqrt2}{2}\)，由上面的分段式知 \(g\left(\dfrac1a\right)>\sqrt2\)，不满足等式. 若 \(-\dfrac{\sqrt2}{2}<a<0\)，则 \(g(a)>\sqrt2\)，而 \(\dfrac1a<-\sqrt2\)，故 \(g\left(\dfrac1a\right)=\sqrt2\)，也不满足等式. 因此 \(a\in\left[-\sqrt2,-\dfrac{\sqrt2}{2}\right]\cup\{1\}\).
\end{enumerate}
\end{solution}

\begin{problem}
设数列 \(\{a_{n}\}\)，\(\{b_{n}\}\)，\(\{c_{n}\}\) 满足：\(b_{n} = a_{n} - a_{n+2}\)，\(c_{n} = a_{n} + 2a_{n+1} + 3a_{n+2}\)（\(n = 1, 2, 3, \cdots\)），证明：\(\{a_{n}\}\) 为等差数列的充分必要条件是 \(\{c_{n}\}\) 为等差数列且 \(b_{n} \leqslant b_{n+1}\)（\(n = 1, 2, 3, \cdots\)）.
\end{problem}

\begin{solution}
设 \(\{c_n\}\) 的公差为 \(d_2\).

若 \(\{a_n\}\) 为等差数列，设其公差为 \(d\)，则 \(a_{n+1}-a_n=d\). 因而 \(b_n=a_n-a_{n+2}=-2d\)，从而 \(b_n=b_{n+1}\). 又
\(c_n=a_n+2(a_n+d)+3(a_n+2d)=6a_n+8d\)，所以 \(c_{n+1}-c_n=6d\) 为常数，\(\{c_n\}\) 为等差数列. 必要性成立.

反之，设 \(\{c_n\}\) 为公差为 \(d_2\) 的等差数列，且 \(b_n\leqslant b_{n+1}\). 由定义得
\(c_n-c_{n+2}=b_n+2b_{n+1}+3b_{n+2}=-2d_2\)，将下标整体加 \(1\) 得
\(b_{n+1}+2b_{n+2}+3b_{n+3}=-2d_2\).
两式相减，得到
\((b_{n+1}-b_n)+2(b_{n+2}-b_{n+1})+3(b_{n+3}-b_{n+2})=0\).
由于三个括号内的量均非负，只能分别为零，故 \(b_n=b_{n+1}\) 对所有 \(n\) 成立. 记这个常数为 \(d_3\)，即 \(a_{n+2}=a_n-d_3\).

于是 \(c_n=4a_n+2a_{n+1}-3d_3\)，从而
\(c_{n+1}-c_n=2(a_{n+1}-a_n)-2d_3=d_2\). 因此 \(a_{n+1}-a_n=\dfrac12d_2+d_3\) 为常数，\(\{a_n\}\) 是等差数列. 充分性成立，命题得证.
\end{solution}
